\chapter{Criptografía RSA}

\section{Introducción}
RSA es uno de los sistemas criptográficos más importantes. Fue desarrollado por Rivest, Shamir y Adleman en 1977.

El sistema RSA utiliza dos claves: una pública $(e,n)$ y una privada $(d,n)$.

\section{Generación de Claves}
El proceso de generación de claves RSA consiste en:
\begin{enumerate}
  \item Elegir dos números primos grandes $p$ y $q$.
  \item Calcular $n = p \cdot q$.
  \item Calcular $\varphi(n) = (p-1)(q-1)$.
  \item Elegir $e$ tal que $1 < e < \varphi(n)$ y $\gcd(e,\varphi(n)) = 1$.
  \item Calcular $d = e^{-1} \bmod \varphi(n)$.
\end{enumerate}

\section{Cifrado y Descifrado}

\subsection{Cifrado}
Para cifrar un mensaje $M$:
\[
C = M^e \bmod n.
\]

\subsection{Descifrado}
Para descifrar un mensaje $C$:
\[
M = C^d \bmod n.
\]

\subsection{Demostración de corrección}
La corrección de RSA se basa en el teorema de Euler:
\[
M^{ed} \equiv M \pmod{n}.
\]

\section{Seguridad de RSA}
La seguridad de RSA se basa en la dificultad de factorizar $n = p \cdot q$.

\section{Implementación en Haskell}

\begin{lstlisting}[caption={RSA en Haskell}, label=lst:rsa-haskell]
-- Algoritmo de Euclides extendido
euclidesExt :: Integer -> Integer -> (Integer, Integer, Integer)
euclidesExt a 0 = (a, 1, 0)
euclidesExt a b =
    let (g, x1, y1) = euclidesExt b (a `mod` b)
    in (g, y1, x1 - (a `div` b) * y1)

-- Inverso modular
inversoMod :: Integer -> Integer -> Maybe Integer
inversoMod a m =
    let (g, x, _) = euclidesExt a m
    in if g == 1 then Just (x `mod` m) else Nothing

-- Potenciación modular
potenciaMod :: Integer -> Integer -> Integer -> Integer
potenciaMod base 0 _ = 1
potenciaMod base exp mod
    | even exp = (potenciaMod base (exp `div` 2) mod) ^ 2 `mod` mod
    | otherwise = (base * potenciaMod base (exp - 1) mod) `mod` mod

-- Verificar si es primo (simple)
esPrimo :: Integer -> Bool
esPrimo n = n > 1 && all (\p -> n `mod` p /= 0) (takeWhile (<= floor (sqrt (fromIntegral n))) [2..])

-- Generar clave RSA
generarRSA :: Integer -> Integer -> Integer -> ((Integer, Integer), (Integer, Integer))
generarRSA p q e =
    let n = p * q
        phi = (p-1) * (q-1)
        d = case inversoMod e phi of
              Just d' -> d'
              Nothing -> error "e no es invertible módulo phi(n)"
    in ((e, n), (d, n))

-- Cifrar RSA
cifrarRSA :: Integer -> (Integer, Integer) -> Integer
cifrarRSA m (e, n) = potenciaMod m e n

-- Descifrar RSA
descifrarRSA :: Integer -> (Integer, Integer) -> Integer
descifrarRSA c (d, n) = potenciaMod c d n

-- Conversión texto-número
textoANumero :: String -> Integer
textoANumero = foldl (\acc c -> acc * 256 + fromIntegral (ord c)) 0

numeroATexto :: Integer -> String
numeroATexto 0 = ""
numeroATexto n = let (q, r) = n `divMod` 256 in numeroATexto q ++ [chr (fromIntegral r)]

-- Ejemplo
ejemploRSA :: IO ()
ejemploRSA = do
    let p = 61
        q = 53
        e = 17
        ((pubE, pubN), (privD, privN)) = generarRSA p q e
        mensaje = 123
        cifrado = cifrarRSA mensaje (pubE, pubN)
        descifrado = descifrarRSA cifrado (privD, privN)
    putStrLn $ "Clave pública: (" ++ show pubE ++ ", " ++ show pubN ++ ")"
    putStrLn $ "Clave privada: (" ++ show privD ++ ", " ++ show privN ++ ")"
    putStrLn $ "Mensaje original: " ++ show mensaje
    putStrLn $ "Mensaje cifrado: " ++ show cifrado
    putStrLn $ "Mensaje descifrado: " ++ show descifrado
\end{lstlisting}

\section{Ejercicios}
\begin{enumerate}
  \item Generar un par de claves RSA con primos pequeños.
  \item Cifrar y descifrar un mensaje.
  \item Implementar en Haskell una función que convierta texto a números para RSA.
\end{enumerate}